\subsection{Nullポインタ}

\subsubsection{MS-DOS時代の「Null pointer assignment」エラー}

\myindex{MS-DOS}
一部の古参は、MS-DOS時代の奇妙なエラーメッセージを思い出すかもしれません：「Null pointer assignment」。
これは何を意味するのでしょうか？

*NIXやWindows OSではアドレス0にメモリを書き込むことはできませんが、メモリ保護が全くなかったため、MS-DOSではそれが可能でした。

\myindex{Turbo C++}
\myindex{Borland C++}
そこで、1990年代初頭の古い私のTurbo C++ 3.0（後にBorland C++に改名）を取り出し、これをコンパイルしてみました：

\begin{lstlisting}[style=customc]
#include <stdio.h>

int main()
{
	int *ptr=NULL;
	*ptr=1234;
	printf ("Now let's read at NULL\n");
	printf ("%d\n", *ptr);
};
\end{lstlisting}

信じがたいことですが、動作します。ただし、終了時にエラーが発生します：

\begin{lstlisting}[caption=古いTurbo C 3.0]
C:\TC30\BIN\1
Now let's read at NULL
1234
Null pointer assignment

C:\TC30\BIN>_
\end{lstlisting}

Borland C++ 3.1の\ac{CRT}のソースコード、ファイル\IT{c0.asm}をより深く掘り下げてみましょう：

\begin{lstlisting}[style=customasmx86]
;       _checknull()    check for null pointer zapping copyright message

...

;       Check for null pointers before exit

__checknull     PROC    DIST
                PUBLIC  __checknull

IF      LDATA  EQ  false
  IFNDEF  __TINY__
                push    si
                push    di
                mov     es, cs:DGROUP@@
                xor     ax, ax
                mov     si, ax
                mov     cx, lgth_CopyRight
ComputeChecksum label   near
                add     al, es:[si]
                adc     ah, 0
                inc     si
                loop    ComputeChecksum
                sub     ax, CheckSum
                jz      @@SumOK
                mov     cx, lgth_NullCheck
                mov     dx, offset DGROUP: NullCheck
                call    ErrorDisplay
@@SumOK:        pop     di
                pop     si
  ENDIF
ENDIF

_DATA           SEGMENT

;       Magic symbol used by the debug info to locate the data segment
                public DATASEG@
DATASEG@        label   byte

;       The CopyRight string must NOT be moved or changed without
;       changing the null pointer check logic

CopyRight       db      4 dup(0)
                db      'Borland C++ - Copyright 1991 Borland Intl.',0
lgth_CopyRight  equ     $ - CopyRight

IF      LDATA  EQ  false
IFNDEF  __TINY__
CheckSum        equ     00D5Ch
NullCheck       db      'Null pointer assignment', 13, 10
lgth_NullCheck  equ     $ - NullCheck
ENDIF
ENDIF

...

\end{lstlisting}

MS-DOSのメモリモデルは本当に奇妙でした（\ref{8086_memory_model}）。レトロコンピューティングやレトロゲームのファンでなければ、詳しく見る価値はおそらくないでしょう。
覚えておかなければならないことの一つは、MS-DOSのメモリセグメント（データセグメントを含む）は、コードやデータが格納されるメモリセグメントですが、「真面目な」\ac{OS}とは異なり、アドレス0から始まることです。

そして、Borland C++の\ac{CRT}では、データセグメントは4つのゼロバイトと著作権文字列「Borland C++ - Copyright 1991 Borland Intl.」で始まります。
4つのゼロバイトとテキスト文字列の整合性は終了時にチェックされ、破損している場合はエラーメッセージが表示されます。

しかし、なぜでしょうか？nullポインタへの書き込みは\CCpp での一般的な間違いであり、*NIXやWindowsでそれを行うと、アプリケーションがクラッシュします。
MS-DOSにはメモリ保護がないため、\ac{CRT}はこれを事後的にチェックし、終了時に警告する必要があります。
このメッセージが表示された場合、プログラムがある時点でアドレス0に書き込んだことを意味します。

我々のプログラムがそれを行いました。そして、1234という数字が正しく読み取られた理由：それが最初の4つのゼロバイトの場所に書き込まれたからです。
終了時にチェックサムが正しくない（数字がそこに残されているため）ため、エラーメッセージが表示されました。

私は正しいでしょうか？
私の仮定をチェックするためにプログラムを書き直しました：

\begin{lstlisting}[style=customc]
#include <stdio.h>

int main()
{
	int *ptr=NULL;
	*ptr=1234;
	printf ("Now let's read at NULL\n");
	printf ("%d\n", *ptr);
	*ptr=0; // psst, cover our tracks!
};
\end{lstlisting}

このプログラムは終了時にエラーメッセージなしで実行されます。

nullポインタ割り当てについて警告する方法はMS-DOSに関連していますが、おそらくメモリ保護や\ac{MMU}のない低コスト\ac{MCU}で今日でも使用できるでしょう。

\subsubsection{なぜ誰かがアドレス0に書き込むのでしょうか？}

しかし、なぜ正常なプログラマがアドレス0に何かを書き込むコードを書くのでしょうか？
それは偶然に行われる可能性があります：例えば、ポインタは新しく割り当てられたメモリブロックに初期化されてから、ポインタを通じてデータを返す何らかの関数に渡される必要があります。

\begin{lstlisting}[style=customc]
int *ptr=NULL;

... we forgot to allocate memory and initialize ptr

strcpy (ptr, buf); // strcpy() terminates silently because MS-DOS has no memory protection
\end{lstlisting}

さらに悪い場合：

\begin{lstlisting}[style=customc]
int *ptr=malloc(1000);

... we forgot to check if memory has been really allocated: this is MS-DOS after all and computers had small amount of RAM,
... and RAM shortage was very common.
... if malloc() returned NULL, the ptr will also be NULL.

strcpy (ptr, buf); // strcpy() terminates silently because MS-DOS has no memory protection
\end{lstlisting}

\subsubsection{\CCpp でのNULL}

C/C++のNULLは、しばしば次のように定義されるマクロにすぎません：

\begin{lstlisting}[style=customc]
#define NULL  ((void*)0)
\end{lstlisting}
（\href{https://github.com/wzhy90/linaro_toolchains/blob/8ff8ae680bac04558d10cc9626e12c4c2f6c1348/arm-cortex_a15-linux-gnueabihf/libc/usr/include/libio.h#L70}{libio.hファイル}）

\IT{void*}は、ポインタであるという事実を反映するデータ型ですが、不明なデータ型（\IT{void}）の値へのポインタです。

NULLは通常、オブジェクトの不在を示すために使用されます。
例えば、単方向リンクリストがあり、各ノードには値（または値へのポインタ）と\IT{next}ポインタがあります。
次のノードがないことを示すために、\IT{next}フィールドに0が格納されます。
他のソリューションはただ悪いだけです。
おそらく、ゼロアドレスでメモリブロックを割り当てる必要がある狂った環境があるかもしれません。次のノードの不在をどのように示すでしょうか？
何らかの\IT{マジックナンバー}？ -1かもしれません？それとも追加ビットを使用しますか？

Wikipediaでこれを見つけることができます：

\begin{framed}
\begin{quotation}
実際、ゼロページの元の優先的使用とは全く逆に、FreeBSD、Linux、Microsoft Windows[2]などの一部の現代オペレーティングシステムは、実際にNULLポインタの使用をトラップするためにゼロページをアクセス不可にしています。
\end{quotation}
\end{framed}
（\url{https://en.wikipedia.org/wiki/Zero_page}）

\subsubsection{関数へのNullポインタ}

関数をそのアドレスで呼び出すことが可能です。
例えば、これをMSVC 2010でコンパイルし、Windows 7で実行します：

\begin{lstlisting}[style=customc]
#include <windows.h>
#include <stdio.h>

int main()
{
	printf ("0x%x\n", &MessageBoxA);
};
\end{lstlisting}

結果は\IT{0x7578feae}で、何度実行しても変わりません。
MessageBoxA関数が存在するuser32.dllが常に同じアドレスでロードされるためです。
また、\ac{ASLR}が有効になっていないためでもあります（その場合、結果は毎回異なるでしょう）。

アドレスで\IT{MessageBoxA()}を呼び出してみましょう：

\begin{lstlisting}[style=customc]
#include <windows.h>
#include <stdio.h>

typedef int (*msgboxtype)(HWND hWnd, LPCTSTR lpText, LPCTSTR lpCaption,  UINT uType);

int main()
{
	msgboxtype msgboxaddr=0x7578feae;

	// force to load DLL into process memory, 
	// since our code doesn't use any function from user32.dll, 
	// and DLL is not imported
	LoadLibrary ("user32.dll");

	msgboxaddr(NULL, "Hello, world!", "hello", MB_OK);
};
\end{lstlisting}

奇妙ですが、Windows 7 x86で動作します。

これはシェルコードで一般的に使用されます。そこから名前でDLL関数を呼び出すのが難しいからです。
そして\ac{ASLR}は対策です。

今、本当に奇妙なことですが、一部の組み込みCプログラマはこのようなコードに馴染みがあるかもしれません：

\begin{lstlisting}[style=customc]
int reset()
{
	void (*foo)(void) = 0;
	foo();
};
\end{lstlisting}

誰がアドレス0で関数を呼び出したいのでしょうか？
これはゼロアドレスにジャンプするポータブルな方法です。
多くの低コストの安価なマイクロコントローラーにもメモリ保護や\ac{MMU}がなく、リセット後、何らかの初期化コードが格納されているアドレス0でコードの実行を開始します。
したがって、アドレス0へのジャンプは自分自身をリセットする方法です。
インラインアセンブリを使用することもできますが、それが不可能な場合、このポータブルメソッドを使用できます。

これは私のLinux x64のGCC 4.8.4でも正しくコンパイルされます：

\begin{lstlisting}[style=customasmx86]
reset:
        sub     rsp, 8
        xor     eax, eax
        call    rax
        add     rsp, 8
        ret
\end{lstlisting}

スタックポインタがシフトされるという事実は問題ではありません：マイクロコントローラーの初期化コードは通常、レジスタと\ac{RAM}の状態を完全に無視し、スクラッチから起動します。

そしてもちろん、このコードはメモリ保護のために*NIXやWindowsでクラッシュし、保護がない場合でも、アドレス0にはコードがありません。

GCCには、そこで関数を呼び出すのではなく、特定のアドレスにジャンプすることを可能にする非標準拡張もあります：
\url{http://gcc.gnu.org/onlinedocs/gcc/Labels-as-Values.html}。